\documentclass[11pt,letterpaper]{article}

\usepackage[T1]{fontenc}
\usepackage{amsmath,amssymb}
\usepackage{tikz}
\usepackage{pgfplots}
\pgfplotsset{compat=1.18}
\usepackage{algorithm}
\usepackage{algpseudocode}
\usepackage{booktabs}
\usepackage{enumitem}
\usepackage[hidelinks]{hyperref}

\newcommand{\N}{\mathbb{N}}
\newcommand{\Z}{\mathbb{Z}}
\newcommand{\Q}{\mathbb{Q}}
\newcommand{\R}{\mathbb{R}}
\newcommand{\C}{\mathbb{C}}
\newcommand{\Map}[3]{#1 \colon #2 \to #3}
\newcommand{\Mod}[1]{#1\text{-}\mathrm{Mod}}
\newcommand{\GroupHom}{\mathrm{Hom}_{\mathbf{Grp}}}
\newcommand{\RingHom}{\mathrm{Hom}_{\mathbf{Ring}}}

\newtheorem{theorem}{Theorem}
\newtheorem{lemma}[theorem]{Lemma}
\newtheorem{corollary}[theorem]{Corollary}
\newtheorem{definition}[theorem]{Definition}
\newtheorem{remark}[theorem]{Remark}
\newenvironment{proof}{\noindent\textit{Proof.}}{\hfill\(\square\)}

\title{The two-affine positivity comparison}
\author{}
\date{\today}

\begin{document}
\maketitle

\begin{abstract}
For an odd integer \(K\ge 3\) and an odd natural number \(q\), the comparison
\(2^K\le\bigl(q+\tfrac13\bigr)^2\) is equivalent to the integer comparison
\(9\cdot 2^K\le(3q+1)^2\). The two sides of the integer comparison are never
equal. Each such exponent has one odd natural candidate in the square bracket
of \(2^K\). The argument uses no square-root extraction in \(\mathbb{R}\) and
does not assert that a cycle exists.
\end{abstract}

\section{The cleared comparison}

Throughout, \(\N\) contains zero. An integer is odd when it has the form
\(2t+1\) for some integer \(t\ge 0\), and even when it has the form \(r+r\).

\begin{definition}
For \(q,K\in\N\), write
\[
\operatorname{TwoAffinePositivity}(q,K)
\]
for the inequality \(9\cdot 2^K\le(3q+1)^2\).
\end{definition}

\begin{lemma}[TwoAffineExpanded]
For every \(q\in\N\),
\[
(3q+1)^2=9q^2+6q+1.
\]
\end{lemma}

\begin{proof}
Expand the square and collect terms.
\end{proof}

\begin{lemma}[SquareModThree]
For every \(q\in\N\),
\[
(3q+1)^2\equiv 1\pmod{3}.
\]
\end{lemma}

\begin{proof}
The expansion gives
\[
(3q+1)^2=3(3q^2+2q)+1.
\]
The first summand is divisible by \(3\), so the residue is \(1\).
\end{proof}

\begin{lemma}[TwoAffineNotEqual]
For all \(q,K\in\N\),
\[
9\cdot 2^K\neq(3q+1)^2.
\]
\end{lemma}

\begin{proof}
The left side equals \(3\cdot(3\cdot 2^K)\), hence is divisible by \(3\).
The right side is congruent to \(1\) modulo \(3\).
\end{proof}

\begin{lemma}[TwoAffinePositivityStrict]
For all \(q,K\in\N\), the inequality
\(9\cdot 2^K\le(3q+1)^2\) holds if and only if
\[
9\cdot 2^K<(3q+1)^2.
\]
\end{lemma}

\begin{proof}
The forward implication is the preceding inequality together with the
impossibility of equality. The reverse implication is the order on \(\N\).
\end{proof}

\begin{theorem}[TwoAffinePositivityIff]
For all \(q,K\in\N\),
\[
2^K\le\Bigl(q+\frac13\Bigr)^2
\quad\text{in }\Q
\]
if and only if \(9\cdot 2^K\le(3q+1)^2\).
\end{theorem}

\begin{proof}
Since \(3\neq 0\) in \(\Q\),
\[
q=\frac{3q}{3},
\]
and therefore
\[
q+\frac13=\frac{3q+1}{3}=\frac{3q+1}{3},
\]
where the numerator on the right is the rational image of the natural number
\(3q+1\). Squaring and using \(3^2=9\) gives
\[
\Bigl(q+\frac13\Bigr)^2=\frac{(3q+1)^2}{9}.
\]
The positive rational \(9\) may be multiplied through an inequality without
changing its direction, so
\[
2^K\le\frac{(3q+1)^2}{9}
\]
if and only if
\[
9\cdot 2^K\le(3q+1)^2.
\]
Both sides of the last comparison are rational images of natural numbers, and
the order embedding of \(\N\) in \(\Q\) identifies it with the stated natural
comparison. The same identity shows that the comparison is strict whenever it
holds.
\end{proof}

\begin{theorem}[TwoAffineWindowComparison]
Let \(K\ge 3\) be odd and let \(q\) be odd. Then
\[
2^K\le\Bigl(q+\frac13\Bigr)^2
\]
if and only if \(\operatorname{TwoAffinePositivity}(q,K)\).
\end{theorem}

\begin{proof}
The hypotheses fix the odd window. The equivalence itself is the denominator
clearing proved for every pair of natural numbers.
\end{proof}

\section{The odd square bracket}

\begin{lemma}[OddPowTwoNotSquare]
For every \(k\in\N\) there is no \(m\in\N\) with
\[
m^2=2^{2k+1}.
\]
\end{lemma}

\begin{proof}
Proceed by induction on \(k\). If \(k=0\), the claim is \(m^2=2\). The cases
\(m=0\) and \(m=1\) give \(0\) and \(1\). If \(m\ge 2\), then
\(m^2\ge 2^2=4>2\).

Suppose the claim holds for \(k\), and suppose \(m^2=2^{2(k+1)+1}=2^{2k+3}\).
The integer \(m\) cannot be odd. An odd integer has the form \(2t+1\), and
\[
(2t+1)^2=4t(t+1)+1
\]
is odd, whereas \(2^{2k+3}\) is even. Thus \(m=t+t=2t\) for some \(t\). Then
\[
(2t)^2=4t^2=2^{2k+3}=4\cdot 2^{2k+1}.
\]
Cancel the positive factor \(4\) to obtain \(t^2=2^{2k+1}\), which the
inductive hypothesis forbids.
\end{proof}

\begin{definition}
For \(K\in\N\), the integer floor \(\operatorname{IntegerSqrtFloor}(K)\) is the
greatest \(m\in\N\) such that \(m^2\le 2^K\). Equivalently,
\[
m^2\le 2^K<(m+1)^2.
\]
The odd candidate \(\operatorname{UniqueOddCandidate}(K)\) is \(m-1\) when
\(m\) is even and is \(m\) when \(m\) is odd.
\end{definition}

\begin{lemma}[IntegerSqrtFloorAtLeastTwo]
If \(K\ge 3\), then
\[
\operatorname{IntegerSqrtFloor}(K)\ge 2.
\]
\end{lemma}

\begin{proof}
The hypothesis gives \(2^K\ge 2^3=8\ge 4=2^2\). The greatest integer whose
square does not exceed \(2^K\) is therefore at least \(2\).
\end{proof}

\begin{lemma}[OddPowStrictSqrt]
Let \(K\ge 3\) be odd and write
\[
m=\operatorname{IntegerSqrtFloor}(K).
\]
Then
\[
m^2<2^K<(m+1)^2.
\]
\end{lemma}

\begin{proof}
The floor inequalities give \(m^2\le 2^K<(m+1)^2\). Equality
\(m^2=2^K\) would mean that an odd power of two is a square, which has been
excluded. The left comparison is therefore strict.
\end{proof}

\begin{lemma}[PredecessorSquareGap]
If \(m\ge 1\), then
\[
\bigl(3(m-1)+1\bigr)^2+4\bigl(3(m-1)+2\bigr)=9m^2.
\]
\end{lemma}

\begin{proof}
Put \(s=m-1\), so \(m=s+1\). The identity becomes
\[
(3s+1)^2+4(3s+2)=9(s+1)^2.
\]
Both sides expand to \(9s^2+18s+9\).
\end{proof}

\begin{lemma}[EvenFloorPositivityFails]
Let \(K\ge 3\) be odd. If
\[
m=\operatorname{IntegerSqrtFloor}(K)
\]
is even, then the odd candidate \(q=m-1\) fails the positivity test at
\((q,K)\).
\end{lemma}

\begin{proof}
The floor bound gives \(m\ge 2\), so \(q\ge 1\). The strict floor inequality
gives \(m^2+1\le 2^K\), and multiplication by \(9\) gives
\[
9(m^2+1)\le 9\cdot 2^K.
\]
The predecessor identity rearranges to
\[
(3q+1)^2+4(3q+2)+9=9m^2+9=9(m^2+1).
\]
Since \(4(3q+2)+9=12q+17\ge 17\),
\[
(3q+1)^2+17\le 9\cdot 2^K.
\]
If \(9\cdot 2^K\le(3q+1)^2\), adding \(17\) would yield
\(9\cdot 2^K+17\le 9\cdot 2^K\), which is false.
\end{proof}

\begin{theorem}[UniqueOddCandidateSpec]
Let \(K\ge 3\) be odd and write \(q=\operatorname{UniqueOddCandidate}(K)\).
Then \(q\) is odd, \(q\ge 1\), and
\[
q^2<2^K<(q+2)^2.
\]
\end{theorem}

\begin{proof}
Let \(m=\operatorname{IntegerSqrtFloor}(K)\), so \(m\ge 2\) and
\(m^2<2^K<(m+1)^2\). Every natural number is even or odd.

If \(m\) is even, then \(q=m-1\). Writing \(m=t+t\) with \(t\ge 1\) gives
\[
q=2(t-1)+1,
\]
so \(q\) is odd and \(q\ge 1\). The successor expansion
\[
m^2=(m-1)^2+2(m-1)+1
\]
gives \((m-1)^2<m^2<2^K\). Also \((m-1)+2=m+1\), so
\[
2^K<(m+1)^2=(q+2)^2.
\]

If \(m\) is odd, then \(q=m\). The floor inequality gives \(q^2<2^K\). Since
\(m+1\le m+2\),
\[
2^K<(m+1)^2\le(m+2)^2=(q+2)^2.
\]
\end{proof}

\begin{theorem}[UniqueOddCandidateUnique]
Let \(K\ge 3\) be odd. Suppose \(r\) is odd and
\[
r^2<2^K<(r+2)^2.
\]
Then \(r=\operatorname{UniqueOddCandidate}(K)\).
\end{theorem}

\begin{proof}
Let \(m=\operatorname{IntegerSqrtFloor}(K)\). From \(r^2<2^K\) one has
\(r^2\le 2^K\), hence \(r\le m\). The integer \(m\) cannot satisfy
\(r+2\le m\): that bound would give
\[
(r+2)^2\le m^2<2^K,
\]
contrary to \(2^K<(r+2)^2\). Thus \(r\le m<r+2\), so \(m=r\) or \(m=r+1\).

If \(m=r\), then \(m\) is odd and the odd candidate is \(m\). If \(m=r+1\),
then \(m\) is the successor of an odd integer, hence even, and the odd
candidate is \(m-1=r\).
\end{proof}

\begin{theorem}[WindowPositivityIff]
Let \(K\ge 3\) be odd and let \(q\) be odd. Suppose
\[
q^2<2^K<(q+2)^2.
\]
Then \(\operatorname{TwoAffinePositivity}(q,K)\) holds if and only if \(q\) is
the odd candidate and the odd candidate itself satisfies the integer test.
\end{theorem}

\begin{proof}
The bracket hypothesis forces \(q\) to be the odd candidate. Substituting that
equality into the integer comparison gives both directions.
\end{proof}

\begin{remark}
The integer test does not assert that any orbit is cyclic. It does not assert
a contraction estimate, and it does not assert the comparison of the two
classical growth constants. A finite set of odd pairs may be checked by exact
arithmetic; such a check is not the theorem. No real square root is required:
the floor is the greatest natural number whose square does not exceed the
power of two, and the positivity test is the integer comparison above.
\end{remark}

\end{document}
